\chapter{Open problems and further mathematical development}\label{app:open}

The records below are the open problems of the registry for this volume. Each is stated with its precise setting, so that progress can be measured against the theorems of the text.

\begin{problem}[Static certification for composite or even \(g\)]\rec{O12.2}\label{op:composite}
Replace hypothesis (S5) of Theorem~\ref{thm:certificate} by a condition on the plateau steps with \(\gcd(g,w)>1\), so that a finite residue test certifies the survivor and the whole future for composite and even \(g\). Known: the survivor exists (Theorem~\ref{thm:dominance}), the class cycle is entered within \(H(g)\) carries (Theorem~\ref{thm:entry}(a)), and all reset states of the six cycles \(g=15,39,51,57,68,87\) fail (S5) (Theorem~\ref{thm:T73}).
\end{problem}

\begin{problem}[Exact or near-sharp transient time]\rec{O12.1}\label{op:exact}
Find the exact transient time, or a bound within a small factor of the observed entry index. Observed entry indices are of order \(10^2\) (the recorded run has a maximum of 208 at \(g=211\); Figure~\ref{fig:bounds} gives the sample computed for this volume), while the proved bounds \(\min(H,K_{\mathrm{poly}})\) exceed \(10^8\) there. The route suggested by the proof of Theorem~\ref{thm:Kpoly} is to bound the epochs by the structure of the saturated odometers rather than by their number of states.
\end{problem}

\begin{problem}[The polynomial bound beyond the carry-level predictor]\rec{O12.3}
Prove a polynomial entry bound for composite \(g\) outside the scope of Theorem~\ref{thm:predictor} (where \(\gcd(g,w)>1\) can occur on the way).
\end{problem}

\begin{problem}[Closed form of the worst-case carry and the survivor for \(q\)-carry families]\rec{O12.4}
Give \(m_p\) and the survivor of the families of Theorem~\ref{thm:qcarry} in closed form from \((g,q)\) alone.
\end{problem}

\begin{problem}[Infinitely many carry sizes]\rec{O12.5}
Decide whether (Q1)--(Q4) of Theorem~\ref{thm:qcarry} are solvable for every odd prime \(q\). Eleven exact instances with \(q\le37\) are recorded. The conditions are finitely many non-membership conditions in a cyclic quotient of \((\Z/g)^\times\), and a Chebotarev or Kummer-type argument appears plausible; it has not been carried out.
\end{problem}

\begin{problem}[The case \(q=11\), \(g=6689\)]\rec{O12.6}
Complete the class-closure certificate for \(q=11\), \(g=6689\); about 50\,000 machine steps are needed.
\end{problem}

\begin{problem}[Realisable triples]\rec{O10.3}
Determine which triples \((g,p,\Abs_p)\) are realised by actual orbits. Theorem~\ref{thm:coprime} and the window \(c^2<g\) restrict them; a complete characterisation is not known.
\end{problem}

\begin{problem}[Composite and even samples]\rec{O10.4}
Theorems~\ref{thm:scheduler} and~\ref{thm:dominance} do not use primality of \(g\), but the numerical sample is almost entirely prime \(g\). A systematic sample of composite and even \(g\) would test the picture of Chapter~\ref{ch:schedule} beyond the six known cycles.
\end{problem}

\section*{Further directions}

\begin{itemize}
\item \emph{Classification of persistent values.} Which \(g\) admit persistent orbits? The answer is known for \(g\mid24\) (none), for \(\gcd(g,6)=1\) (all \(g\ge5\)), and in the instances \(g=15,39,51,57,87,68\) (Remark~\ref{rem:T41}).
\item \emph{Scope of the polynomial bound.} The proof of Theorem~\ref{thm:Kpoly} counts epochs of the saturated odometers; a refinement that exploits the periodicity of \(\Abs_p\) (Theorem~\ref{thm:coprime}) may reduce \(K_{\mathrm{poly}}\) substantially.
\item \emph{Other collective grammars.} The theory of this volume concerns the selected rule \(C\). The same programme (a coordinate in which the dynamics becomes finite, a scheduler, a certificate and effective bounds) can be asked for other collective rules; Volume~I shows that the interaction axioms considered there do not single out one.
\end{itemize}
